\documentclass[12pt,a4paper,openany,oneside]{book}

% --- Paquetes ---
\usepackage[utf8]{inputenc}
\usepackage[spanish, es-noshorthands]{babel}
\usepackage{amsmath, amsfonts, amssymb}
\usepackage{graphicx}
\usepackage{geometry}
\usepackage{hyperref}
\usepackage{booktabs}
\usepackage{fancyhdr}
\usepackage{listings}
\usepackage{xcolor}
\usepackage{tikz}
\usepackage{pgfplots}
\usepackage{colortbl}
\usepackage{multirow}
\usepackage{hhline}
\usetikzlibrary{arrows.meta, positioning, shapes.geometric, calc}
\pgfplotsset{compat=1.18}

\geometry{margin=2.5cm}

% --- Colores y estilos para código Python ---
\definecolor{codegreen}{rgb}{0,0.6,0}
\definecolor{codegray}{rgb}{0.5,0.5,0.5}
\definecolor{codepurple}{rgb}{0.58,0,0.82}
\definecolor{backcolour}{rgb}{0.95,0.95,0.92}

\lstset{
    language=Python,
    backgroundcolor=\color{backcolour},   
    commentstyle=\color{codegreen},
    keywordstyle=\color{blue},
    numberstyle=\tiny\color{codegray},
    stringstyle=\color{codepurple},
    basicstyle=\ttfamily\small,
    breaklines=true,
    frame=single,
    showstringspaces=false
}

% --- Estilo de cabecera y pie ---
\pagestyle{fancy}
\fancyhf{}
\renewcommand{\headrulewidth}{0.4pt}
\renewcommand{\footrulewidth}{0.4pt}
\fancyhead[L]{\small Carlos César Sánchez Coronel}
\fancyhead[R]{\small Sesión 3: Regresión Lineal y Regularización}
\fancyfoot[L]{\small Carlos César Sánchez Coronel}
\fancyfoot[C]{\thepage}

% --- Portada ---
\title{
    \textbf{Sesión 3} \\ 
    \Huge \textbf{REGRESIÓN LINEAL Y REGULARIZACIÓN} \\
    \Large Modelos lineales, interpretación y control de la complejidad
}
\author{
    \small Por \\ \vspace{0.2cm}
    \Large \textsc{Carlos César Sánchez Coronel}
}
\date{2026}

\begin{document}

\maketitle
\tableofcontents
\addcontentsline{toc}{chapter}{Índice general}

% ------------------------------------------------------------------
% CAPÍTULO 3: SESIÓN 3 - REGRESIÓN LINEAL Y REGULARIZACIÓN
% ------------------------------------------------------------------
\chapter{Regresión Lineal y Regularización}

En esta sesión se aborda el modelo de regresión lineal en sus formas simple y múltiple, su interpretación, los supuestos que lo sustentan y las métricas para evaluarlo. Se introduce el problema del sobreajuste y las técnicas de regularización (Ridge, Lasso, Elastic Net) como mecanismos para controlar la complejidad y mejorar la generalización.

\section{Logro de la sesión}
Construir e interpretar modelos de regresión lineal, aplicando regularización para mejorar su generalización, y seleccionar la métrica de evaluación adecuada según el contexto del problema.

\section{Introducción: el modelo lineal como base fundamental}
La regresión lineal es el punto de partida en el modelado predictivo. Su simplicidad, interpretabilidad y base matemática sólida la convierten en una herramienta esencial, incluso cuando se utilizan modelos más complejos. Comprenderla en profundidad permite entender conceptos clave como la función de pérdida, el gradiente descendente, el sesgo-varianza y la regularización, que son comunes a prácticamente todos los algoritmos de machine learning.

\section{Regresión lineal simple}

\subsection{Formulación matemática}
El modelo de regresión lineal simple relaciona una variable independiente \(x\) con una variable dependiente \(y\) mediante una ecuación lineal:
\[
y = \beta_0 + \beta_1 x + \varepsilon
\]
donde:
\begin{itemize}
    \item \(\beta_0\) es la intersección (ordenada al origen).
    \item \(\beta_1\) es la pendiente (coeficiente de la variable).
    \item \(\varepsilon\) es el término de error, que captura la variabilidad no explicada por el modelo.
\end{itemize}

\subsection{Estimación de parámetros: mínimos cuadrados ordinarios (MCO)}
Dado un conjunto de \(n\) pares \((x_i, y_i)\), se buscan los valores \(\hat{\beta}_0\) y \(\hat{\beta}_1\) que minimicen la suma de cuadrados de los residuos (RSS):
\[
RSS = \sum_{i=1}^n (y_i - \hat{\beta}_0 - \hat{\beta}_1 x_i)^2
\]
Las soluciones se obtienen derivando e igualando a cero:
\[
\hat{\beta}_1 = \frac{\sum_{i=1}^n (x_i - \bar{x})(y_i - \bar{y})}{\sum_{i=1}^n (x_i - \bar{x})^2}, \quad
\hat{\beta}_0 = \bar{y} - \hat{\beta}_1 \bar{x}
\]
donde \(\bar{x}\) e \(\bar{y}\) son las medias muestrales.

\subsection{Interpretación de coeficientes}
\begin{itemize}
    \item \(\beta_1\) indica el cambio esperado en \(y\) por cada incremento unitario en \(x\).
    \item \(\beta_0\) es el valor esperado de \(y\) cuando \(x=0\). A menudo no tiene interpretación práctica si \(x=0\) está fuera del rango de los datos.
\end{itemize}

\subsection{Caso de uso: relación entre años de experiencia y salario}
Se tiene una muestra de empleados con sus años de experiencia y salario anual. El modelo permite estimar cuánto aumenta el salario por cada año adicional de experiencia. Por ejemplo, si \(\hat{\beta}_1 = 1500\), cada año adicional se asocia con un incremento promedio de 1500 unidades monetarias.

\section{Regresión lineal múltiple}

\subsection{Formulación matricial}
En presencia de \(p\) variables predictoras, el modelo se escribe como:
\[
y_i = \beta_0 + \beta_1 x_{i1} + \beta_2 x_{i2} + \dots + \beta_p x_{ip} + \varepsilon_i
\]
En forma matricial, para todas las observaciones:
\[
\mathbf{y} = \mathbf{X}\boldsymbol{\beta} + \boldsymbol{\varepsilon}
\]
donde \(\mathbf{y}\) es un vector \(n \times 1\), \(\mathbf{X}\) es la matriz de diseño \(n \times (p+1)\) (incluyendo una columna de unos para el intercepto), \(\boldsymbol{\beta}\) es el vector de coeficientes \((p+1) \times 1\) y \(\boldsymbol{\varepsilon}\) es el vector de errores.

\subsection{Estimación por mínimos cuadrados}
La función de costo a minimizar es:
\[
J(\boldsymbol{\beta}) = (\mathbf{y} - \mathbf{X}\boldsymbol{\beta})^T (\mathbf{y} - \mathbf{X}\boldsymbol{\beta})
\]
La solución analítica (ecuación normal) es:
\[
\hat{\boldsymbol{\beta}} = (\mathbf{X}^T \mathbf{X})^{-1} \mathbf{X}^T \mathbf{y}
\]
Esta solución existe siempre que \(\mathbf{X}^T \mathbf{X}\) sea invertible, lo que requiere que las columnas de \(\mathbf{X}\) sean linealmente independientes (no multicolinealidad perfecta).

\subsection{Interpretación de coeficientes en regresión múltiple}
Cada coeficiente \(\beta_j\) representa el cambio esperado en \(y\) por cada incremento unitario en \(x_j\), manteniendo constantes las demás variables (ceteris paribus). Esta interpretación es clave en análisis de negocio: por ejemplo, en un modelo de precios de viviendas, el coeficiente de "metros cuadrados" indica cuánto aumenta el precio por cada metro cuadrado adicional, asumiendo misma ubicación, número de habitaciones, etc.

\subsection{Caso de uso: predicción de precios de viviendas}
Variables predictoras: área construida, número de habitaciones, antigüedad, distancia al centro. El modelo permite cuantificar el impacto de cada característica en el precio, útil para tasaciones y decisiones de inversión.

\section{Supuestos del modelo de regresión lineal}
Para que la inferencia estadística (intervalos de confianza, tests de hipótesis) sea válida, deben cumplirse ciertos supuestos sobre los errores \(\varepsilon_i\):

\subsection{Linealidad}
La relación entre cada predictor y la respuesta es lineal. Se verifica mediante gráficos de residuos vs. valores ajustados; si aparecen patrones, puede requerirse transformación de variables.

\subsection{Independencia}
Los errores son independientes entre sí. En datos temporales o espaciales, puede haber autocorrelación. Se evalúa con gráficos de residuos en orden de recolección.

\subsection{Homocedasticidad}
La varianza de los errores es constante a lo largo de todos los niveles de los predictores. Si la varianza cambia (heterocedasticidad), los errores estándar pueden estar sesgados. Se detecta con gráficos de residuos vs. ajustados o tests estadísticos (Breusch-Pagan).

\subsection{Normalidad de los residuos}
Los errores siguen una distribución normal. Necesario para tests de hipótesis e intervalos de confianza, especialmente en muestras pequeñas. Se evalúa con histogramas, Q-Q plots o tests de Shapiro-Wilk.

En la práctica, el incumplimiento moderado de estos supuestos no invalida el modelo predictivo, pero sí afecta la interpretación de la significancia estadística. Para fines predictivos, la homocedasticidad y la normalidad son menos críticas.

\section{Métricas de evaluación en regresión}

\subsection{Error Cuadrático Medio (MSE)}
\[
MSE = \frac{1}{n} \sum_{i=1}^n (y_i - \hat{y}_i)^2
\]
Mide el promedio de los errores al cuadrado. Penaliza más los errores grandes debido al cuadrado. Está en unidades al cuadrado de la variable objetivo, lo que dificulta la interpretación directa.

\subsection{Raíz del Error Cuadrático Medio (RMSE)}
\[
RMSE = \sqrt{\frac{1}{n} \sum_{i=1}^n (y_i - \hat{y}_i)^2}
\]
Devuelve el error a las unidades originales. Es la métrica más común y directamente interpretable. Un RMSE de 5000 unidades monetarias significa que, en promedio, el error de predicción es de 5000.

\subsection{Error Absoluto Medio (MAE)}
\[
MAE = \frac{1}{n} \sum_{i=1}^n |y_i - \hat{y}_i|
\]
Promedio de los errores absolutos. Es robusto a outliers, ya que no eleva al cuadrado. Útil cuando no se quiere penalizar excesivamente los errores grandes.

\subsection{Comparación entre RMSE y MAE con ejemplos simulados}
Supóngase un conjunto de errores: [10, 10, 10, 100].
\begin{itemize}
    \item MAE = (10+10+10+100)/4 = 32.5
    \item RMSE = sqrt((100+100+100+10000)/4) = sqrt(2600) ≈ 51.0
\end{itemize}
El RMSE es mayor porque el error de 100 se amplifica al cuadrado. Por tanto, si los errores grandes son especialmente indeseables (ej. en predicción de demanda de productos perecederos, donde un error grande puede significar pérdida de producto), el RMSE es más apropiado. Si todos los errores importan por igual, el MAE es más interpretable.

\subsection{Coeficiente de determinación \(R^2\)}
\[
R^2 = 1 - \frac{\sum_{i=1}^n (y_i - \hat{y}_i)^2}{\sum_{i=1}^n (y_i - \bar{y})^2} = 1 - \frac{RSS}{TSS}
\]
Interpretación: proporción de la varianza total de \(y\) explicada por el modelo. Varía entre 0 y 1 (puede ser negativo si el modelo es peor que la media). Un \(R^2\) de 0.8 indica que el 80\% de la variabilidad es explicada por los predictores.

\subsection{Cuándo usar cada métrica}
\begin{itemize}
    \item \textbf{RMSE:} Cuando los errores grandes son críticos y se desea interpretación en unidades originales. Predicción de precios, demanda, etc.
    \item \textbf{MAE:} Cuando se busca robustez frente a outliers, o en problemas donde el error absoluto tiene sentido de negocio (ej. error en inventarios).
    \item \textbf{MSE:} Útil como función de pérdida en optimización, por su diferencialidad.
    \item \textbf{\(R^2\):} Para evaluar la bondad de ajuste relativa y comparar modelos, pero no sirve para seleccionar entre modelos con diferente número de predictores (se prefiere \(R^2\) ajustado).
\end{itemize}

\section{El problema del sobreajuste y la necesidad de regularización}
Cuando se añaden muchos predictores, el modelo puede ajustarse demasiado a los datos de entrenamiento, capturando ruido en lugar de la señal subyacente. Esto se manifiesta en un alto rendimiento en entrenamiento pero pobre en prueba. La regularización introduce una penalización por la magnitud de los coeficientes, restringiendo la complejidad del modelo.

\subsection{Comprensión intuitiva: sesgo-varianza}
La regularización incrementa ligeramente el sesgo (el modelo se fuerza a ser más simple) pero reduce drásticamente la varianza (menos sensible a las fluctuaciones de los datos). El objetivo es encontrar el equilibrio que minimice el error total en datos no vistos.

\section{Regularización Ridge (L2)}
\subsection{Formulación matemática}
Ridge añade una penalización proporcional a la suma de los cuadrados de los coeficientes (excepto el intercepto):
\[
J_{\text{ridge}}(\boldsymbol{\beta}) = \sum_{i=1}^n (y_i - \beta_0 - \sum_{j=1}^p \beta_j x_{ij})^2 + \lambda \sum_{j=1}^p \beta_j^2
\]
En forma matricial:
\[
J_{\text{ridge}}(\boldsymbol{\beta}) = (\mathbf{y} - \mathbf{X}\boldsymbol{\beta})^T (\mathbf{y} - \mathbf{X}\boldsymbol{\beta}) + \lambda \|\boldsymbol{\beta}\|_2^2
\]
El hiperparámetro \(\lambda \geq 0\) controla la fuerza de la regularización. Si \(\lambda = 0\), se obtiene MCO; si \(\lambda \to \infty\), los coeficientes tienden a cero.

\subsection{Solución analítica}
Derivando e igualando a cero, se obtiene:
\[
\hat{\boldsymbol{\beta}}_{\text{ridge}} = (\mathbf{X}^T \mathbf{X} + \lambda \mathbf{I})^{-1} \mathbf{X}^T \mathbf{y}
\]
donde \(\mathbf{I}\) es la matriz identidad de tamaño \(p\) (se omite el intercepto si no se regulariza). La adición de \(\lambda \mathbf{I}\) asegura que la matriz sea invertible incluso si \(\mathbf{X}^T \mathbf{X}\) es singular (multicolinealidad perfecta).

\subsection{Interpretación y efecto}
\begin{itemize}
    \item Ridge encoge los coeficientes hacia cero, pero ninguno llega exactamente a cero (a menos que \(\lambda\) sea infinito). Por tanto, todos los predictores se mantienen en el modelo.
    \item Es especialmente útil cuando hay multicolinealidad, ya que estabiliza las estimaciones.
    \item La penalización L2 favorece que los coeficientes sean pequeños pero no nulos, lo que mantiene la interpretabilidad (todos los predictores contribuyen).
\end{itemize}

\subsection{Caso de uso: predicción de gasto en clientes con muchas variables correlacionadas}
En marketing, se tienen decenas de variables demográficas y de comportamiento, muchas correlacionadas. Ridge permite incluir todas sin que la multicolinealidad desestabilice las estimaciones.

\section{Regularización Lasso (L1)}
\subsection{Formulación matemática}
Lasso (Least Absolute Shrinkage and Selection Operator) utiliza una penalización L1:
\[
J_{\text{lasso}}(\boldsymbol{\beta}) = \sum_{i=1}^n (y_i - \beta_0 - \sum_{j=1}^p \beta_j x_{ij})^2 + \lambda \sum_{j=1}^p |\beta_j|
\]

\subsection{Característica clave: selección de variables}
La penalización L1 tiene la propiedad de que, para valores suficientemente grandes de \(\lambda\), algunos coeficientes se vuelven exactamente cero. Esto realiza una selección automática de variables, produciendo modelos más simples e interpretables.

\subsection{Solución computacional}
No existe solución analítica cerrada; se utilizan algoritmos de optimización como el descenso por coordenadas o LARS. A diferencia de Ridge, Lasso puede producir modelos dispersos (sparse).

\subsection{Caso de uso: identificación de factores clave en un proceso industrial}
Se tienen cientos de sensores en una planta. Lasso puede seleccionar aquellos pocos cuyas lecturas son realmente relevantes para predecir la calidad del producto, facilitando el diagnóstico y la interpretación.

\section{Elastic Net}
\subsection{Formulación matemática}
Elastic Net combina las penalizaciones L1 y L2:
\[
J_{\text{elastic}}(\boldsymbol{\beta}) = \frac{1}{2n} \sum_{i=1}^n (y_i - \beta_0 - \sum_{j=1}^p \beta_j x_{ij})^2 + \lambda \left( \alpha \sum_{j=1}^p |\beta_j| + \frac{1-\alpha}{2} \sum_{j=1}^p \beta_j^2 \right)
\]
El parámetro \(\alpha \in [0,1]\) controla la mezcla: \(\alpha = 0\) es Ridge, \(\alpha = 1\) es Lasso. Valores intermedios combinan ambas propiedades.

\subsection{Ventajas respecto a Lasso}
\begin{itemize}
    \item Lasso tiene limitaciones cuando hay grupos de variables correlacionadas: tiende a seleccionar solo una del grupo. Elastic Net puede seleccionar grupos completos.
    \item Cuando \(p > n\) (más predictores que observaciones), Lasso puede seleccionar a lo sumo \(n\) variables; Elastic Net no tiene esa restricción.
\end{itemize}

\subsection{Caso de uso: análisis de expresión génica}
En bioinformática, se tienen miles de genes (\(p\)) y pocas muestras (\(n\)). Elastic Net permite identificar grupos de genes relacionados con una enfermedad, manteniendo la estabilidad de Ridge y la selección de Lasso.

\section{Selección del hiperparámetro \(\lambda\) (y \(\alpha\))}
El valor óptimo de \(\lambda\) se determina mediante validación cruzada. Se evalúa el error (generalmente MSE) para distintos \(\lambda\) en las particiones de validación y se elige el \(\lambda\) que minimiza el error promedio.

\begin{itemize}
    \item \textbf{Validación cruzada k-fold:} Dividir datos en \(k\) partes, entrenar con \(k-1\) y validar con la restante, repetir \(k\) veces.
    \item \textbf{Curva de regularización:} Representar coeficientes en función de \(\lambda\).
\end{itemize}

\section{Interpretación de coeficientes en modelos regularizados}
En modelos regularizados, los coeficientes ya no son insesgados (tienen un sesgo introducido por la penalización). Por tanto, su interpretación como "efecto ceteris paribus" es menos precisa, pero siguen siendo útiles para ranking de importancia. En Lasso, los coeficientes exactamente cero indican variables excluidas.

\section{Comparación entre Ridge, Lasso y Elastic Net}

\begin{table}[h]
\centering
\begin{tabular}{|l|l|l|l|}
\hline
\textbf{Característica} & \textbf{Ridge (L2)} & \textbf{Lasso (L1)} & \textbf{Elastic Net} \\
\hline
Selección de variables & No & Sí & Sí (en grupos) \\
Manejo de multicolinealidad & Excelente & Regular & Bueno \\
Estabilidad & Alta & Puede ser inestable & Intermedia \\
Interpretabilidad & Media (todos los coeficientes) & Alta (modelo disperso) & Alta \\
Uso típico & Muchos predictores correlacionados & Selección de variables & Grupos de variables correlacionadas \\
\hline
\end{tabular}
\caption{Comparación de métodos de regularización.}
\end{table}

\section{Caso de estudio integrado: Predicción de precios de viviendas con regularización}

\subsection{Datos}
Dataset con 80 variables predictoras (área, habitaciones, año, ubicación, calidad, etc.) y precios de venta.

\subsection{Preprocesamiento}
\begin{itemize}
    \item Imputación de nulos (mediana para numéricas, moda para categóricas).
    \item Codificación one-hot de variables categóricas (aumenta dimensionalidad a >200).
    \item Estandarización de todas las numéricas (necesaria para regularización, ya que las penalizaciones son sensibles a la escala).
\end{itemize}

\subsection{Modelado}
Se prueban tres enfoques:
\begin{enumerate}
    \item Regresión lineal sin regularización (MCO).
    \item Ridge con validación cruzada para elegir \(\lambda\).
    \item Lasso con validación cruzada.
    \item Elastic Net con búsqueda de \(\lambda\) y \(\alpha\).
\end{enumerate}

\subsection{Resultados esperados}
\begin{itemize}
    \item MCO puede tener coeficientes inestables y alto error en prueba por overfitting.
    \item Ridge estabiliza las estimaciones y mejora el error de prueba.
    \item Lasso produce un modelo más simple, seleccionando quizás 20-30 variables relevantes, con rendimiento similar a Ridge.
    \item Elastic Net puede equilibrar ambas ventajas.
\end{itemize}

\subsection{Evaluación}
Se comparan RMSE y MAE en el conjunto de prueba. Además, se analiza la interpretabilidad: los coeficientes de Lasso y Elastic Net permiten identificar qué características son más determinantes en el precio (ubicación, área, calidad), mientras que Ridge incluye todas pero con coeficientes pequeños.

\section{Resumen y conexión con el resto del curso}
La regresión lineal, en sus versiones simple y múltiple, proporciona la base para entender conceptos clave de modelado. La regularización introduce la noción de control de complejidad, esencial para evitar sobreajuste. Estos conceptos se extenderán a otros modelos:
\begin{itemize}
    \item En regresión logística, la regularización evita overfitting en clasificación.
    \item En redes neuronales, la regularización L2 (weight decay) es estándar.
    \item La selección de variables con Lasso anticipa técnicas de selección de características en modelos más complejos.
\end{itemize}
La comprensión de las métricas y su elección según el contexto de negocio es una habilidad transversal que se aplicará en todos los modelos futuros.

\end{document}